\documentclass[11pt]{article}
\usepackage[margin=1in]{geometry}
\usepackage{amsmath}
\usepackage{amssymb}
\usepackage{booktabs}
\usepackage{algorithm}
\usepackage{algpseudocode}
\usepackage[colorlinks=true,urlcolor=blue]{hyperref}

\newcommand{\x}{\mathbf{x}}
\newcommand{\onehot}{\mathrm{onehot}}

\title{Unifying reward-tilted posterior guidance with D-CBG\\
\large What proposal does D-CBG induce over $\x_0$?}
\author{}
\date{2026-08-01}

\begin{document}
\maketitle

\begin{abstract}
D-CBG and our reward-tilted posterior guidance are usually described as
competing classifier-based mechanisms, but they act at different levels of the
generative process: D-CBG tilts the one-step posterior over $\x_{t-1}$, while
our method tilts the denoiser's distribution over the clean sequence $\x_0$.
This note reads off exactly what D-CBG computes from the implementation, lifts
its classifier factor to an $\x_0$-level proposal, and shows that the resulting
importance weights collapse to a single clean expression in which the
classifier's normalizer cancels. The combined estimator uses the classifier to
decide \emph{where the $N$ atoms land} and the exact reward to decide
\emph{what they are worth}, and remains consistent for the reward-tilted target
regardless of how bad the classifier is.
\end{abstract}

\section{Notation}

Let $L$ be the sequence length, $V$ the vocabulary, $c$ the target class (in
our experiments, ``QED above the 90th percentile of QM9''), and
$p_\phi(c \mid \x, t)$ the \emph{noisy} classifier D-CBG trains on diffusion
latents. Write $\x_t^{l \leftarrow k}$ for the sequence $\x_t$ with position $l$
overwritten by token $k$. Throughout, define the \textbf{substitution score}
\begin{equation}
s_l(k) \;:=\; \log p_\phi\!\left(c \;\middle|\; \x_t^{l \leftarrow k},\, t\right)
\in (-\infty, 0].
\label{eq:score}
\end{equation}
This is exactly the tensor called \texttt{classifier\_log\_prob} in
\texttt{diffusion.py}, of shape $(B, L, V)$.

\section{What D-CBG actually computes}

Reading \texttt{Diffusion.\_cbg\_denoise} (\texttt{diffusion.py:1344--1456}),
D-CBG performs a \emph{per-position} tilt of the unguided reverse posterior:
\begin{equation}
\boxed{\;
q_{\text{CBG}}\!\left(\x_{t-1}^l = k \mid \x_t\right)
  \;\propto\; q\!\left(\x_{t-1}^l = k \mid \x_t\right)\,
              \exp\!\left(\gamma\, s_l(k)\right)
\;}
\label{eq:dcbg}
\end{equation}
implemented as
\texttt{guided\_log\_probs = gamma * classifier\_log\_prob + diffusion\_log\_probs}
followed by a softmax over the vocabulary axis, with the positions then sampled
independently. For absorbing-state diffusion the unmasked positions are forced
to copy $\x_t$, so \eqref{eq:dcbg} is only active on masked positions.

The two \texttt{use\_approx} branches differ only in how $s_l(k)$ is obtained:
\begin{itemize}
  \item \textbf{Exact} (\texttt{use\_approx=False}): materialize all $L \cdot V$
        single-token substitutions $\x_t^{l \leftarrow k}$ and run the
        classifier on each. Cost $L \cdot V$ classifier forwards per step ---
        this is the variant that ran out of memory at batch size 64.
  \item \textbf{Approximate} (\texttt{use\_approx=True}): a first-order Taylor
        expansion in the one-hot input,
        \begin{equation}
        s_l(k) \approx \log p_\phi(c \mid \x_t, t)
          + \frac{\partial \log p_\phi(c \mid \x, t)}{\partial x_{l,k}}
            \bigg|_{\x = \x_t}
          - \frac{\partial \log p_\phi(c \mid \x, t)}{\partial x_{l, \x_t^l}}
            \bigg|_{\x = \x_t},
        \end{equation}
        the last term centering the score at the token currently occupying
        position $l$. Cost: one forward and one backward per step.
\end{itemize}

\paragraph{Three properties worth naming.}
First, D-CBG's tilt lives at the level of $\x_{t-1}$, not $\x_0$. Second, the
tilt is \emph{factorized}: it is a product of per-position factors, so it can
never express a preference that couples positions. Third, D-CBG applies the
tilt and stops --- there is no correction step, so the sampler targets
\eqref{eq:dcbg} rather than any exactly-specified distribution over $\x_0$. The
approximation error of $p_\phi$ (and of the Taylor expansion) propagates
straight into the samples.

\section{What our method computes}

Our estimator tilts at the level of the clean sequence. With
$\x_0^{(n)} \sim p_\theta(\x_0 \mid \x_t)$ drawn i.i.d.\ from the denoiser,
\begin{equation}
w_n = \frac{\exp(\lambda\, r(\x_0^{(n)}))}
           {\sum_{m} \exp(\lambda\, r(\x_0^{(m)}))},
\qquad
\bar{\x}_0 = \sum_{n=1}^{N} w_n \, \onehot(\x_0^{(n)}),
\end{equation}
and $\x_{t-1}$ is drawn from the posterior with $\bar{\x}_0$ substituted for
$\x_0$. This is self-normalized importance sampling for the target
$q(\x_0) \propto p_\theta(\x_0)\exp(\lambda r(\x_0))$, so unlike \eqref{eq:dcbg}
it is \emph{consistent}: as $N \to \infty$ it recovers the exact reward-tilted
distribution. The reward $r$ sees clean $\x_0$, so it may couple positions
arbitrarily --- QED is a whole-molecule functional --- and needs no training.

The price is the atomic support discussed in the companion note: only the $N$
sampled candidates carry mass, and they are sampled from the \emph{untilted}
denoiser, so the atoms land wherever the base model happens to look rather than
where the reward is high.

\section{The unification}

The two mechanisms are complementary precisely because of the above:
D-CBG has a cheap factorized signal that can move probability mass but cannot
be corrected, and we have an exact uncorrectable-by-nature signal that can only
reweight mass that is already there. The combination is to use the first as the
\textbf{proposal} and the second as the \textbf{weight}.

\subsection{What the optimal factorized proposal is}
\label{sec:optimal}

The obstacle is that \eqref{eq:dcbg} tilts $\x_{t-1}$, whereas importance
sampling needs a distribution over $\x_0$. Before borrowing anything from
D-CBG, it is worth deriving what we \emph{should} be looking for, because that
derivation is what tells us the classifier is the right object rather than a
convenient one.

Two facts frame the search.

\paragraph{(i) Any factorized proposal is legal.} Self-normalized importance
sampling is consistent for the target
$q(\x_0 \mid \x_t) \propto p_\theta(\x_0 \mid \x_t) e^{\lambda r(\x_0)}$ for
\emph{any} proposal $\tilde{p}$ with
$\mathrm{supp}(\tilde{p}) \supseteq \mathrm{supp}(q)$, provided the weights
carry the ratio $q/\tilde{p}$. So the choice of proposal is not a correctness
question at all --- it cannot make the estimator wrong, only slow. What it
controls is the effective sample size, i.e.\ how many of the $N$ atoms actually
count. We restrict to \emph{factorized} $\tilde{p}$ purely for tractability: a
factorized distribution is the only kind we can draw $N$ i.i.d.\ samples from at
the cost of one categorical draw.

\paragraph{(ii) The best factorized proposal is the target's marginals.} Among
factorized distributions $\tilde{p}(\x_0) = \prod_l \tilde{p}_l(\x_0^l)$, the
one minimizing $\mathrm{KL}(q \,\|\, \tilde{p})$ is obtained by matching
marginals. Writing it out,
\begin{equation}
\mathrm{KL}(q \,\|\, \textstyle\prod_l \tilde{p}_l)
 = -H(q) - \sum_{l}\sum_{k} q_l(k) \log \tilde{p}_l(k),
\end{equation}
which is maximized term by term at $\tilde{p}_l = q_l$, the per-position
marginal of the target. So the object we want is
\begin{equation}
q_l(k) \;=\; \sum_{\x_0 : \, \x_0^l = k} q(\x_0 \mid \x_t)
      \;\propto\; p_\theta\!\left(\x_0^l = k \mid \x_t\right)
       \cdot \underbrace{\mathbb{E}_{\x_0^{-l} \sim p_\theta(\cdot \mid \x_t)}
         \!\left[\, e^{\lambda r(\x_0^{-l},\, k)} \,\right]}_{
         \textstyle =:\; g_l(k)},
\label{eq:ideal}
\end{equation}
where the factorization of $p_\theta$ was used to drop the conditioning inside
the expectation. Equation~\eqref{eq:ideal} is the derivation that was missing:
\textbf{the exact per-position multiplier is not an arbitrary
$\exp(\cdot)$ --- it is the expected exponentiated reward obtained by pinning
position $l$ to token $k$ and letting the model fill in the rest.}

$g_l(k)$ is intractable: evaluating it by Monte Carlo needs $L \cdot V$ reward
calls per sample, and $r$ is a black box, so there is no closed form. Any usable
method must replace $g_l$ with an estimate.

\subsection{The classifier is an amortized estimate of \texorpdfstring{$g_l$}{g\_l}}

Now the reason D-CBG's score belongs here. Take the tilt in its
class-conditional form, $q(\x_0) \propto p_\theta(\x_0)\, p(c \mid \x_0)^\gamma$,
which is what D-CBG's classifier is trained for ($c$ = ``QED above the 90th
percentile''). Substituting into \eqref{eq:ideal} with $\gamma = 1$,
\begin{equation}
g_l(k)
 = \mathbb{E}_{\x_0^{-l} \sim p_\theta(\cdot \mid \x_t)}
   \!\left[\, p\!\left(c \mid \x_0^{-l}, k\right) \right]
 = \Pr\!\left(c \;\middle|\; \x_t,\; \x_0^l = k\right).
\label{eq:gl-as-prob}
\end{equation}
That is exactly the quantity a noisy classifier is trained to produce: the
probability of the target class given the current latent, with one position
pinned. D-CBG's score $s_l(k) = \log p_\phi(c \mid \x_t^{l \leftarrow k}, t)$ is
a \emph{learned amortization of \eqref{eq:gl-as-prob}} --- it replaces an
$L \cdot V$-fold Monte Carlo integral over the model's own completions with one
network evaluation. That, and not analogy with \eqref{eq:dcbg}, is why
multiplying the denoiser by $e^{\gamma s_l(k)}$ is the natural move: it is the
plug-in estimator of the optimal factorized proposal \eqref{eq:ideal}.

Accordingly, define
\begin{equation}
\tilde{p}_\theta\!\left(\x_0^l = k \mid \x_t\right)
  \;=\; \frac{1}{Z_l}\,
        p_\theta\!\left(\x_0^l = k \mid \x_t\right)\,
        \exp\!\left(\gamma\, s_l(k)\right),
\qquad
Z_l = \sum_{k'} p_\theta\!\left(\x_0^l = k' \mid \x_t\right)
      e^{\gamma s_l(k')} .
\label{eq:proposal}
\end{equation}
It is normalizable in closed form, factorized (so drawing $N$ candidates still
costs one categorical sample), and costs one classifier evaluation per step
shared across all $N$ candidates.

\paragraph{Where the plug-in is inexact, and why that is tolerable.}
Three gaps separate $e^{\gamma s_l(k)}$ from $g_l(k)$:
\begin{enumerate}
  \item \textbf{Substitution versus conditioning.} $p_\phi$ is evaluated at
        $\x_t^{l \leftarrow k}$ --- token $k$ written into the \emph{noisy}
        latent --- whereas \eqref{eq:gl-as-prob} conditions on
        $\x_0^l = k$ in the \emph{clean} sequence. These coincide only in the
        low-noise limit.
  \item \textbf{Binary class versus continuous reward.} $p_\phi$ predicts a
        thresholded indicator, while our target uses $e^{\lambda r}$ with $r$
        continuous. The exponent $\gamma$ is a free parameter that partially
        absorbs this mismatch, which is another reason it must be retuned
        rather than copied from D-CBG.
  \item \textbf{Estimation error} in $p_\phi$ itself.
\end{enumerate}
By fact (i) none of these threaten correctness: the weights in
\eqref{eq:final} divide the entire factor back out, so a wrong $g_l$ estimate
costs effective sample size and nothing else. This is the whole reason the
combination is worth doing --- it lets us use a crude, cheap, possibly
misspecified estimate of an intractable quantity without inheriting its bias.

\paragraph{Even a perfect $g_l$ leaves work for the weights.} Note that
matching the marginals exactly would \emph{not} make the weights uniform. The
target $q$ does not factorize --- $r$ is a whole-molecule functional, and that
coupling is precisely the part a per-position proposal cannot represent. What
remains in the weights after an ideal proposal is exactly the position-coupling
information, which is the part only the exact reward can supply. The division of
labour is therefore not merely convenient but structural: the classifier can in
principle supply all of the factorizable signal and none of the rest.

\paragraph{A classifier-free alternative.} Equation~\eqref{eq:ideal} also
suggests a variant needing no classifier at all: estimate $g_l(k)$ from the
candidates we already draw, by taking the $M \ll N$ highest-reward candidates
and re-scoring single-position substitutions on them. This costs $M \cdot L
\cdot V$ reward calls per step, which the optimized reward path makes borderline
affordable, and it removes the $t$-mismatch of item 1 entirely because the
reward is evaluated on clean sequences. Recorded as an option; the classifier
route is far cheaper.

\subsection{The importance weights}

Sampling from $\tilde{p}_\theta$ instead of $p_\theta$ requires the weights to
carry the likelihood ratio. For the target
$q(\x_0) \propto p_\theta(\x_0)\exp(\lambda r(\x_0))$,
\begin{align}
w_n
&\;\propto\; \exp\!\left(\lambda r(\x_0^{(n)})\right)
   \frac{p_\theta(\x_0^{(n)} \mid \x_t)}{\tilde{p}_\theta(\x_0^{(n)} \mid \x_t)}
\label{eq:w1}\\[2pt]
&\;=\; \exp\!\left(\lambda r(\x_0^{(n)})\right)
   \prod_{l=1}^{L} Z_l \,
   \exp\!\left(-\gamma\, s_l\!\left(\x_0^{(n),l}\right)\right)
\label{eq:w2}\\[2pt]
&\;\propto\; \exp\!\left(
   \lambda\, r(\x_0^{(n)})
   \;-\; \gamma \sum_{l=1}^{L} s_l\!\left(\x_0^{(n),l}\right)
 \right),
\label{eq:w3}
\end{align}
where the last step drops $\prod_l Z_l$ because it does not depend on $n$ and
therefore cancels in the self-normalization. So the whole combination reduces
to a single extra term in the log-weight:
\begin{equation}
\boxed{\;
\log w_n \;=\; \underbrace{\lambda\, r(\x_0^{(n)})}_{\text{exact reward}}
  \;-\; \underbrace{\gamma \sum_{l} \log p_\phi\!\left(c \,\middle|\,
        \x_t^{\,l \leftarrow \x_0^{(n),l}}\right)}_{\text{undo the proposal tilt}}
  \;+\; \text{const},
\;}
\label{eq:final}
\end{equation}
followed by a softmax over $n$. Nothing else in the algorithm changes:
$\bar{\x}_0$, the posterior substitution, and the time window all stay as they
are.

\subsection{Why this is the right object}

Read \eqref{eq:proposal}--\eqref{eq:final} as a statement about roles.
Section~\ref{sec:optimal} established that the best factorized proposal is the
target's per-position marginals \eqref{eq:ideal}, and that D-CBG's classifier is
an amortized estimate of exactly that object. In this framing,
\begin{itemize}
  \item \textbf{D-CBG alone} uses the surrogate and never corrects it, so it is
        biased with respect to the true tilt by however wrong $p_\phi$ is;
  \item \textbf{ours alone} uses the exact tilt but places its atoms with the
        untilted denoiser, so it is consistent but support-limited;
  \item \textbf{combined}, the surrogate places the atoms where the reward is
        plausibly high and the exact reward corrects whatever the surrogate got
        wrong.
\end{itemize}
Crucially the correction makes the combination \emph{robust to a bad
classifier}: $p_\phi$ enters only through the proposal, and \eqref{eq:final}
divides its influence back out. A poor classifier costs effective sample size,
never correctness. That is a strictly better contract than D-CBG's, where
classifier error is indistinguishable from signal.

\paragraph{Limiting cases.} Both reduce correctly.
\begin{itemize}
  \item $\gamma \to 0$: the proposal becomes $p_\theta$ and the correction term
        vanishes, recovering our current method exactly.
  \item $\lambda \to 0$: the weights become
        $w_n \propto \prod_l e^{-\gamma s_l(\x_0^{(n),l})}$, which exactly
        undoes the proposal tilt, so $\bar{\x}_0 \to \x_\theta$ and sampling
        reduces to the unguided base model. This is a useful implementation
        check: with $\lambda = 0$ the combined sampler must reproduce unguided
        statistics no matter what $\gamma$ is.
\end{itemize}

\subsection{A bias/variance knob}

Equation \eqref{eq:final} fully removes the classifier tilt. Removing it only
partially,
\begin{equation}
\log w_n = \lambda\, r(\x_0^{(n)})
  - \beta\gamma \sum_l s_l\!\left(\x_0^{(n),l}\right),
\qquad \beta \in [0, 1],
\label{eq:beta}
\end{equation}
interpolates between two \emph{different targets} rather than merely trading
variance:
\begin{itemize}
  \item $\beta = 1$ targets $q(\x_0) \propto p_\theta(\x_0) e^{\lambda r(\x_0)}$
        --- the reward-tilted distribution, classifier used purely as a
        proposal.
  \item $\beta = 0$ targets the product of experts
        $q(\x_0) \propto p_\theta(\x_0)\, e^{\lambda r(\x_0)}
        \prod_l p_\phi(c \mid \x_t^{l \leftarrow \x_0^l})^{\gamma}$
        --- the classifier's opinion is kept as an additional factor in the
        objective.
\end{itemize}
Neither is obviously preferable: $\beta = 0$ is not a bug if one is willing to
define the target that way, and it will have higher ESS because the weights no
longer fight the proposal. $\beta = 1$ is the principled choice if the reward is
the definition of what we want.

\subsection{Scale warning}

The two terms in \eqref{eq:final} are not on comparable scales, and reusing
D-CBG's tuned $\gamma$ directly will not work. In D-CBG, $\gamma s_l(k)$
competes against a \emph{per-position} log-posterior. In \eqref{eq:final} it
appears as $\gamma \sum_l s_l$, a sum over all $L = 32$ positions, competing
against the scalar $\lambda r$. Expect the useful $\gamma$ here to be smaller
than D-CBG's by roughly a factor of $L$, and expect the correction term to
dominate the reward unless $\gamma$ is retuned from scratch. Monitoring ESS
during the sweep is the practical guard: a collapse to $\mathrm{ESS} \approx 1$
means the correction term is swamping the reward.

\section{Algorithm}

\begin{algorithm}[H]
\caption{One guided denoising step (combined)}
\begin{algorithmic}[1]
\State $\log \x_\theta \gets$ denoiser$(\x_t, \sigma_t)$
       \Comment{1 network forward}
\State $s_l(k) \gets \log p_\phi(c \mid \x_t^{l \leftarrow k}, t)$ for all $l, k$
       \Comment{1 forward + backward (approx.), or $LV$ forwards (exact)}
\State $\tilde{p}_\theta(\x_0^l{=}k \mid \x_t)
        \propto \x_\theta[l, k]\, e^{\gamma s_l(k)}$
       \Comment{normalize per position}
\State $\x_0^{(n)} \sim \tilde{p}_\theta(\cdot \mid \x_t)$, $n = 1 \dots N$
       \Comment{one categorical draw; free in $N$}
\State $\log \tilde{w}_n \gets \lambda\, r(\x_0^{(n)})
        - \beta\gamma \sum_l s_l(\x_0^{(n),l})$
       \Comment{$N$ RDKit calls, deduplicated and pooled}
\State $w \gets \mathrm{softmax}_n(\log \tilde{w})$;\quad
       $\bar{\x}_0 \gets \sum_n w_n \onehot(\x_0^{(n)})$
\State $\x_{t-1} \sim q(\x_{t-1} \mid \x_t, \bar{\x}_0)$
\end{algorithmic}
\end{algorithm}

\paragraph{Cost.} One network forward plus one classifier forward/backward per
step, independent of $N$; the reward calls scale with the number of
\emph{distinct} candidates. Relative to our current sampler this adds only the
classifier pass, so with \texttt{use\_approx=True} the combined method should
cost roughly what D-CBG costs, while retaining our method's $N$-scaling. Use the
approximate branch: the exact one needs $L \cdot V$ classifier forwards and
already exhausted a 3090's memory at batch size 64.

\section{Open issues}

\begin{enumerate}
  \item \textbf{How good an estimate of $g_l$ is $s_l$?} This is now a
        well-posed empirical question rather than a vague worry:
        \eqref{eq:gl-as-prob} says the target is
        $\Pr(c \mid \x_t, \x_0^l = k)$, and it can be estimated directly, if
        expensively, by Monte Carlo over model completions. Measuring the
        agreement between $e^{\gamma s_l}$ and that estimate as a function of
        $t$ would say exactly where the classifier route is worth using. The
        time-window results predict poor agreement at large $t$, which argues
        for switching the classifier tilt on only inside the same window where
        the reward is informative.
  \item \textbf{Weight blow-up.} The correction $e^{-\gamma \sum_l s_l}$ is
        large exactly for candidates the classifier disliked. If the classifier
        and the reward disagree, a single candidate can dominate. Clipping the
        log-weights, or $\beta < 1$, are the obvious mitigations.
  \item \textbf{Absorbing state.} On unmasked positions $\x_{t-1}$ is forced to
        copy $\x_t$, so both the proposal tilt and the sum in \eqref{eq:final}
        should run over masked positions only; including copied positions adds
        a constant to every $\log w_n$, which cancels, but only if the same set
        of positions is used for all $n$ --- which it is, since masking depends
        on $\x_t$ alone.
  \item \textbf{Does it actually help?} The empirical claim to test is that at
        matched $N$ the combined method attains higher ESS and higher novel QED
        than either component. The companion note's finding that our method's
        gain generalizes to novel molecules while D-CBG's concentrates on
        recall suggests the combination should be evaluated on novel QED
        specifically, not on the all-valid mean.
\end{enumerate}

\end{document}
